
abéliens, de la
façon suivante : si @@M0@@, on associe à @@M1@@ l'ensemble des séries formelles @@M2@@, vérifiant les conditions suivantes : si @@M3@@, alors
@@M4@@
Cet ensemble sera noté @@M5@@, ou @@M6@@ si aucune confusion n'est possible ; il est clair que c'est un groupe abélien pour la multiplication des séries formelles, et qu'on définit bien ainsi un foncteur
@@M7@@
Si @@M8@@ est muni de la filtration grossière, @@M9@@ est le groupe des séries formelles d'augmentation @@M10@@, qu'on notera aussi @@M11@@.

\textbf{Lemme 1.2.} \uline{Le foncteur} @@M12@@ \uline{est représentable}.

Soit @@M13@@, une famille d'indéterminées ; on attribue le poids @@M14@@ à @@M15@@, et on filtre l'anneau @@M16@@ par le poids par rapport aux @@M17@@. On définit alors une bijection~:
@@M18@@
en associant à toute série @@M19@@ l'homomorphisme @@M20@@ déterminé par @@M21@@. Cet homomorphisme est compatible aux filtrations, et on obtient ainsi un isomorphisme de foncteurs.

On pose @@M22@@~; pour toute @@M23@@-algèbre filtrée @@M24@@ et toute série @@M25@@, il existe un homomorphisme et un
seul 
